\section{Comparison Systems}
\label{sec:comparison}

Every system in \cref{tab:comparison} is trained on the same Track 1 training file, selects its settings on nothing but training data, and is scored once on the same 660 development instances by the official scorer, run unchanged. Where a system holds out validation data, it uses the same grouped 20\% split by source pair (\cref{sec:protocol}). The settings of the existing systems are read from their recorded \path{config.json} files, and \cref{app:comparison} gives them in full.

\begin{table*}[t]
\centering
\small
\begin{tabular}{@{}p{0.10\textwidth}p{0.19\textwidth}p{0.42\textwidth}p{0.20\textwidth}@{}}
\toprule
Group & System & Representation and training & Settings recorded in \\
\midrule
3.1 Simple & majority & constant \FW{} (tie with \BW{} broken by label order) & \path{trivial-baselines} \\
 & uniform random & one uniform label per instance, seed 13 & same \\
 & always \EQ{} & constant & same \\
 & always \NO{} & constant & same \\
\midrule
3.2 Classical ML & cues + logreg & 17 standardised surface cues; logistic regression, $C = 1$, balanced class weights & \path{ml-baseline} \\
 & tfidf + logreg & word 1--2 gram TF-IDF of ``$a$ \texttt{[SEP]} $b$'', sublinear tf; same regression & same \\
 & both + logreg & cues and TF-IDF concatenated; same regression & same \\
 & cues + gbdt & 17 cues; histogram gradient boosting, 300 iterations, learning rate 0.05, balanced class weights & same \\
\midrule
3.3 Pretrained & DeBERTa-v3-base & fully fine-tuned, linear classification head, 5 epochs, 3 seeds & \path{deberta-base} \\
 & Qwen2.5-1.5B, zero-shot & frozen, prompted; scores of the four option letters & plan (\cref{sec:approach}) \\
 & Qwen2.5-1.5B, LoRA $\lambda = 0$ & frozen backbone, LoRA adapters, both orders, cross-entropy only; single-pass decoding & plan (\cref{sec:approach}) \\
\midrule
3.4 Proposed & Qwen2.5-1.5B, LoRA $\lambda > 0$ & as above plus the consistency term of \cref{eq:loss}; symmetrized decoding (\cref{eq:sym}) & plan (\cref{sec:approach}) \\
\bottomrule
\end{tabular}
\caption{Comparison systems. The folders are \texttt{experiments/2026-09-29-<name>-track1/}. All systems share the training data, the grouped validation split where one is used, the development set and the official scorer.}
\label{tab:comparison}
\end{table*}

\paragraph{What each group tells us.}
The simple baselines mark the floor. Always \EQ{} is the instructive one: \SoftCons{} 1.000 at weighted F1 0.110, so \SoftCons{} is never read on its own. The classical baselines measure how far surface form goes. The cue model reaches \SoftCons{} 0.794 largely because its features behave predictably under reversal, which takes no understanding of the pair. DeBERTa-v3-base \citep{he-etal-2023-debertav3}, fully fine-tuned with a linear classification head, is the bar to reach: weighted F1 0.807, \SoftCons{} 0.877 and \HardCons{} 0.810 as the mean of three seeds. The zero-shot decoder shows what the backbone knows before adaptation.

\paragraph{The comparison that isolates our contribution.}
The plain-LoRA decoder ($\lambda = 0$) has the same prompt, adapters, data, both training orders and schedule as the proposed system. The two differ only in the components we propose: the Jensen--Shannon consistency term in the training loss and symmetrized decoding at prediction time. Each can be switched off independently (\cref{sec:ablations}).
